\section{An all-rank bound for odd-characteristic symplectic Weil cores}
\label{sec:symplectic-weil-core}

\paragraph{Status and scope.}
The following new author derivation has passed an independent whole-proof,
source and arithmetic review by \texttt{/root/g34\_eightfolds}.
The root reviewer has separately passed the complete mathematical
argument and finite arithmetic; its fresh source review is recorded
separately.  The author's finite arithmetic certificate has been executed; the
unbounded-rank assertions are proved by the displayed inequalities,
not by that finite check.  No assertion about a non-Weil irreducible
representation of a symplectic group is made.  In particular this
section does not remove the general small-field Lie-core gate.

For a finite-order complex operator $T$ on a space of dimension $a$, set
\[
 \nu(T)=a-\max_{\lambda}\dim\ker(T-\lambda).
\]
This quantity is unchanged by multiplying $T$ by a scalar.  It is
bounded above by the ordinary moving rank $\operatorname{rank}(T-1)$.
We allow arbitrary scalar phases throughout.

\subsection*{D-SPW-01. The Heisenberg model and its exact trace identity}

Let $q=p^f$ be odd, $n\ge2$, and let $V$ be a $2n$-dimensional
symplectic space over $\mathbf F_q$ with form $\langle\ ,\ \rangle$.
The Heisenberg group is
\[
 H(V)=V\times\mathbf F_q,\qquad
 (v,z)(w,t)=(v+w,z+t+\tfrac12\langle v,w\rangle).
\]
Fix a nontrivial additive character $\psi$ of $\mathbf F_q$.
Choose a Lagrangian decomposition of $V$ and realize the irreducible
Heisenberg representation with central character $\psi$ on the
space $\mathcal W$ of functions on one Lagrangian.  Its dimension is
$Q=q^n$.  Write $D(v)$ for the operator of $(v,0)$.

The linear Weil representation and its Schr\"odinger model are
published inputs, specifically G\'erardin
\cite[Theorem~2.4(a), pp.~63--64, proof pp.~64--69]{SPW:Gerardin}.
The model on the unipotent radical of a Lagrangian stabilizer is
in equations (2.8)--(2.9) on p.~65.  The following exact trace
and irreducibility statements are also published there; we use
them directly rather than reprove them.

\begin{lemma}[D-SPW-01: published trace input and scalar normalization]
\label{D-SPW-01}
For the full Weil representation $U$ of a finite symplectic
space in odd characteristic,
\[
          |\operatorname{tr}U(A)|^2=|\ker(A-1)|.            \tag{SPW.1}
\]
For the space $V$ above, the parity operator
$P:F(x)\mapsto F(-x)$ has two irreducible nonisomorphic
$\operatorname{Sp}(V)$ eigenspaces, of dimensions
$(Q+1)/2$ and $(Q-1)/2$.
The same absolute-trace formula holds for any unitary
Heisenberg intertwiner implementing $A$, and that intertwiner
commutes with parity.  This last formulation permits the
trace-form specialization used for field automorphisms below.
\end{lemma}

\begin{proof}
Formula (SPW.1) is exactly G\'erardin
\cite[Theorem~4.4(c), pp.~84--85]{SPW:Gerardin}.
The parity decomposition and its irreducibility are
\cite[Corollary~4.4(a) and its proof, p.~85]{SPW:Gerardin}.
These published assertions hold over every odd finite field,
so apply both to $\mathbf F_q$ and to $\mathbf F_p$ later.

It remains only to justify the formulation for an arbitrary
unitary intertwiner.  Irreducibility of the Heisenberg
representation makes two intertwiners for the same $A$ scalar
multiples of one another.  If both are unitary the scalar has
absolute value one, so (SPW.1) is unchanged.
The operators $P$ and $U_A P U_A^{-1}$ implement the same
central transformation $-1$ of the finite symplectic space,
so they too differ by a scalar.  Both have trace one, since
zero is the unique fixed point of $x\mapsto-x$ in odd
characteristic.  That scalar is one, as required.
\end{proof}

We will also use the uniqueness in Theorem~2.4(a): on a fixed
irreducible Heisenberg space any two linear Weil extensions
differ by a one-dimensional character of the symplectic group.
For $n\ge2$ the group is perfect, so that character is trivial.
This identifies the two realizations in the field-restriction
argument below without an additional extension theorem.

\subsection*{D-SPW-02. Every scalar and outer normalizer}

\begin{theorem}[D-SPW-02: a uniform one-fifth bound]
\label{D-SPW-02}
Let $\rho$ be any complex irreducible Weil constituent of
$\operatorname{Sp}_{2n}(q)$, with $q$ odd and $n\ge2$, and let
$N=\rho(\operatorname{Sp}_{2n}(q))$.  Put
$a=\dim\rho=(q^n\pm1)/2$.
Every nonscalar finite-order element $T$ of
$N_{\operatorname{GL}_a(\mathbf C)}(N)$ satisfies
\[
                              \nu(T)\ge a/5.               \tag{SPW.2}
\]
This includes every compatible field automorphism and arbitrary
scalar phases.  The same conclusion applies to any faithful
quasisimple core representation whose projective representation
is one of these Weil representations.
\end{theorem}

\begin{proof}
We first bound inner elements.  If $g\in\operatorname{Sp}_{2n}(q)$
is noncentral, put
\[
 d_+=\dim\ker(g-1),\qquad d_-=\dim\ker(g+1)
\]
on the natural $2n$-dimensional space.  These spaces are disjoint,
so $d_++d_-\le2n$.  If either dimension is $2n-1$, the other is
zero.  To see this, the generalized eigenspace of a symplectic
matrix for the eigenvalue $1$, and likewise for $-1$, is
nondegenerate alternating: different primary spaces are paired
only with their reciprocal eigenvalues.  Its dimension is
therefore even.  An eigenspace of dimension $2n-1$ consequently
has generalized eigenspace of dimension $2n$, leaving no other
eigenvalue.

If neither dimension is $2n-1$, both are at most $2n-2$.
Convexity of $x\mapsto q^{x/2}$ on the interval with fixed sum
at most $2n$ gives
\[
 q^{d_+/2}+q^{d_-/2}
 \le\max\{q^{n-1/2}+1,\ q^{n-1}+q\}
 =q^{n-1/2}+1.                                             \tag{SPW.3}
\]
For the final equality, after factoring $\sqrt q-1$ it is
enough that $q^{n-1}\ge\sqrt q+1$, which holds for $q\ge3$
and $n\ge2$.

The trace on a parity constituent is
$(\operatorname{tr}U_g\pm\operatorname{tr}(PU_g))/2$.
Since $PU_g$ implements $-g$, (SPW.1) and (SPW.3) imply
\[
 \frac{|\operatorname{tr}\rho(g)|}{a}
 \le \frac{q^{n-1/2}+1}{q^n-1}
 =q^{-1/2}+\frac{1+q^{-1/2}}{q^n-1}.                       \tag{SPW.4}
\]
There is no assumption that the largest eigenspace has eigenvalue
one.  For any unitary operator of dimension $a$, the contribution
of an eigenspace of maximal dimension $a-\nu$ gives
\[
        |\operatorname{tr}T|\ge a-2\nu(T).                \tag{SPW.5}
\]
Thus a trace ratio at most $3/5$ implies (SPW.2).

For $q\ge5$, (SPW.4) is less than
$1/2+(3/2)/24=9/16<3/5$.  For $q=3,n\ge4$, its maximum
is at $n=4$, where
\[
 \frac{27\sqrt3+1}{80}
 <\frac{239}{400}<\frac35,
 \qquad \sqrt3<\frac{26}{15}.                              \tag{SPW.6}
\]
At $q=3,n=3$, the same bound is less than $2/3$.
The two degrees are $13,14$, so (SPW.5) and integrality give
$\nu\ge3\ge a/5$.  At $q=3,n=2$, the degrees are $4,5$,
and the elementary nonscalar inequality $\nu\ge1$ suffices.
This proves (SPW.2) for all inner elements.

We now cover the actual outer normalizers.  The published
automorphism input is that for odd $q$ all automorphisms of
$\operatorname{Sp}_{2n}(q)$ are products of inner, diagonal and
field automorphisms; there is no graph automorphism.
See Taylor \cite[\S1.5, author p.~2; \S12.5, author p.~24]
{SPW:Taylor}.  The latter passage gives the full automorphism
group, not merely a subgroup of known automorphisms.  The
universal cover of $\operatorname{PSp}_{2n}(q)$ is
$\operatorname{Sp}_{2n}(q)$ in the range at hand.  In particular
$\operatorname{PSp}_4(3)$ has no exceptional multiplier in its
characteristic-three description; see
\cite[Chapter~5, Table~1 and preceding paragraph, numbered p.~35]
{BI:HHM}.  Consequently every automorphism of a represented
central quotient lifts to this universal perfect cover.

The diagonal automorphism with nonsquare similitude multiplier
does not preserve any Weil constituent.  Here is a direct check,
including the scalar issue.  In the Schr\"odinger model a
transvection with parameter $b\ne0$ acts as
\[
                  F(x)\longmapsto\psi(bx_1^2/2)F(x).
\]
Its full trace is $q^{n-1}G_\psi(b)$, where
$G_\psi(b)=\sum_{y\in\mathbf F_q}\psi(by^2/2)$, and its
parity-twisted trace is one.  Multiplying $b$ by a nonsquare
changes the nonzero Gauss sum to its negative.  Indeed
$G_\psi(b)=\sum_y\eta(y)\psi(by/2)$ for the quadratic
character $\eta$, and substitution gives
$G_\psi(cb)=\eta(c)G_\psi(b)$; its squared absolute value is
$q$ by the elementary orthogonality calculation.
Thus the traces $(q^{n-1}G_\psi(b)\pm1)/2$ on each fixed
parity differ after this automorphism.  The two representations
cannot even be projectively equivalent: a projective equivalence
between linear representations of the perfect group would
differ by a one-dimensional character, which is trivial.

For completeness, we construct and bound field intertwiners.
Write
\[
 \psi(x)=\exp\left(\frac{2\pi i}{p}
                 \operatorname{Tr}_{\mathbf F_q/\mathbf F_p}(a x)\right),
 \qquad a\in\mathbf F_q^*.
\]
Regard $V$ as a $2nf$-dimensional $\mathbf F_p$ symplectic
space with alternating form
$B_a(v,w)=\operatorname{Tr}_{\mathbf F_q/\mathbf F_p}
(a\langle v,w\rangle)$.
This form is nondegenerate by nondegeneracy of the field trace
and of $\langle\ ,\ \rangle$.
The Heisenberg representation for this $\mathbf F_p$ space is
the same representation on the same $q^n$-dimensional function
space: the center of the $\mathbf F_q$ Heisenberg group acts
through its trace character.  The uniqueness just recalled
identifies the models.  Restricting the linear Weil representation
of $\operatorname{Sp}_{2nf}(p)$ to $\operatorname{Sp}_{2n}(q)$
gives the chosen full Weil representation, because both extend
this Heisenberg representation and the latter group is perfect.
This proves the needed restriction assertion rather than
assuming it as an additional functoriality theorem.

Let $\sigma$ be coordinatewise $p$-power Frobenius and set
\[
            t=a^{(p-1)/2},\qquad S=t\sigma.
\]
Then $t^2=a^p/a$, and
\[
 B_a(Sv,Sw)=\operatorname{Tr}(a^p\langle v,w\rangle^p)
          =B_a(v,w).
\]
Hence $S\in\operatorname{Sp}_{2nf}(p)$, and conjugation by $S$
on $\operatorname{Sp}_{2n}(q)$ is the field automorphism
$g\mapsto g^{(p)}$.  It has finite order, also directly from
$S^f=\operatorname{Nm}_{q/p}(t)I=\pm I$.
The Weil operator of $S$ commutes with parity by D-SPW-01,
so every field automorphism preserves each Weil constituent.
Together with the diagonal calculation, this shows that a
normalizing matrix can induce only an inner-field automorphism.
A diagonal-field product cannot restore a constituent that the
diagonal factor exchanges.

Consider a nontrivial field component, represented by the
semilinear map $A=gS^k$ with $0<k<f$ and
$g\in\operatorname{Sp}_{2n}(q)$.
Put $d=\gcd(f,k)$, so $d\le f/2$.
The fixed vectors of $A$ have dimension at most $2n$ over
$\mathbf F_{p^d}$.  To verify this bound, suppose fixed vectors
independent over that fixed field have a minimal nonzero linear
relation over $\mathbf F_q$.  Normalize one coefficient to one,
apply $A$, and subtract the original relation.  Minimality forces
every coefficient to be fixed by $\sigma^k$.  The original
relation is then over $\mathbf F_{p^d}$, a contradiction.
The same argument applies to $-A$.  Their fixed spaces therefore
have $\mathbf F_p$ dimensions at most $2nd$.
Using (SPW.1) over $\mathbf F_p$ and the parity trace formula gives
\[
 \frac{|\operatorname{tr}(U_A|\mathcal W^{\pm})|}{a}
 \le\frac{2p^{nd}}{p^{nf}-1}
 \le\frac{2q^{n/2}}{q^n-1}
 \le\frac{18}{80}<\frac35.                                \tag{SPW.7}
\]
The last inequality uses $q\ge9$, $n\ge2$ and the decreasing
function $2x/(x^2-1)$ for $x\ge9$.
Every other linear intertwiner for this automorphism differs
by a scalar, by irreducibility.  Thus (SPW.5) proves (SPW.2)
for every field coset as well.

Finally, the normalizer modulo scalars embeds in the automorphism
group of the simple projective image: the linear centralizer is
scalar, and a central automorphism of a perfect group is trivial.
The projective kernel of a nontrivial Weil constituent is exactly
the center, since the projective quotient is simple.  The preceding
lift and intertwiner arguments therefore cover all normalizers
of the represented image, including those of central quotients.
A faithful quasisimple representation with the same projective
representation has the same image after the unique lift from
the universal perfect cover; any remaining difference is a scalar
character and is trivial.  This proves the final assertion.
\end{proof}

\subsection*{D-SPW-03. Consequence for actual sixfold Hurwitz sources}

\begin{theorem}[D-SPW-03: symplectic Weil cores reduce to degrees four and five]
\label{D-SPW-03}
Suppose an inducing block of a finite-cover sixfold source has
an irreducibly acting normal central product of quasisimple
groups, each represented by an odd-characteristic symplectic
Weil constituent with $n\ge2$.  Allow all compatible scalar
and outer automorphisms, and arbitrary permutations of the
tensor factors.  If the actual period map has rank nine, then
there is only one factor and
\[
                 (q,n,a)=(3,2,4)\quad\hbox{or}\quad(3,2,5).
                                                               \tag{SPW.8}
\]
In particular every factor with $q\ge5$ or $n\ge3$, and every
tensor product with two or more factors, is impossible.
This assertion does not require a nonsplit target polarization
or equality of the inducing and original character fields.
\end{theorem}

\begin{proof}
Let $e$ be the total inducing degree.  For a factorwise tensor
operator, fixing all eigenvalues in the other factors shows
that its largest eigenspace has dimension at most the product
of their dimensions times the largest eigenspace of any chosen
nonscalar factor.  Thus D-SPW-02 gives $\nu\ge e/5$.
If an operator permutes tensor factors nontrivially, the weighted
cyclic-word proof of D-SPOR-03 gives $\nu\ge e/4$.
That proof only uses equal dimensions along a permutation cycle,
finite order and diagonalizability of the closing twist; it is
independent of the nature of the quasisimple factors.  Hence
every nonscalar element in the full normalizer satisfies
$\nu\ge e/5$.  A nonidentity scalar inertia has moving rank $e$,
so the same lower bound holds for every genuine inertia.

The actual intermediate-cover factorization established in
D-BI-18 applies even when the inducing character field is larger.
Rank nine forces the intermediate base to have genus zero and
at least twelve genuine branch points, while the exact
Chevalley--Weil Euler identity is
\[
             \sum_i\operatorname{rank}(g_i-1)=2e+6.
\]
Consequently
\[
                   12\lceil e/5\rceil\le2e+6.             \tag{SPW.9}
\]
It follows first that $e\le15$, and then exact integer division
leaves, for $e\ge4$, only
\[
                              e\in\{4,5,9,10,15\}.         \tag{SPW.10}
\]
Every factor has degree at least four.  Two factors already
give $e\ge16$, which is impossible.

For the single remaining factor, its degree is $(q^n\pm1)/2$.
The inequality $e\le15$ implies $q^n\le31$.
For odd prime powers $q$ and $n\ge2$, the only possibilities
are $(q,n)=(3,2),(3,3),(5,2)$.
The latter two have degrees $13,14$ and $12,13$, respectively,
none in (SPW.10).  The first gives precisely (SPW.8).
All of these are necessary restrictions on an actual period
map; no dimension equality is used to assert dominance.
\end{proof}

\paragraph{Finite check and remaining cases.}
The separate arithmetic script
\texttt{lie\_small\_fields/check\_symplectic\_weil.py} checks
the rational inequalities in (SPW.6)--(SPW.7), the exceptional
degree ceilings, and the complete finite degree list after the
proved cutoff $e\le15$.  Its actual output is retained.
It is an author certificate, not an independent peer audit.
The surviving single-factor degree-four and degree-five cases
are not asserted impossible by this section; they belong to the
separate low-dimensional inducing analyses.  Non-Weil
representations, characteristic-two symplectic groups, and
other small-field Lie families require separate arguments.
